\begin{table}[!htbp]
\centering
\caption{Desempenho dos seis métodos por coleção.}
\label{tab:res-desempenho}
\begin{singlespace}
\normalsize
\setlength{\tabcolsep}{5pt}
\renewcommand{\arraystretch}{1.15}
\begin{tabular}{lrrrrrr}
\hline\hline
\textbf{Método} & \textbf{\shortstack{\textit{Gap} 55\\(\%)}} & \textbf{\shortstack{\textit{Gap} 19\\(\%)}} & \textbf{\shortstack{\textit{Gap} 74\\(\%)}} & \textbf{\shortstack{Posição\\média}} & \textbf{\shortstack{Melhor\\média}} & \textbf{\shortstack{Melhor média\\exclusiva}} \\
\hline
IBS & 2,65 & 5,92 & 3,49 & 2,07 & 44 & 37 \\
MIP & 18,52 & 5,04 & 15,06 & 2,16 & 11 & 6 \\
CUTH & 29,39 & 3,23 & 22,67 & 2,80 & 20 & 20 \\
LBBD & 38,11 & 9,67 & 30,81 & 4,39 & 11 & 4 \\
ILS & 39,57 & 10,03 & 31,98 & 5,32 & 1 & 0 \\
RS & 43,75 & 9,05 & 34,84 & 4,26 & 0 & 0 \\
\hline\hline
\end{tabular}
\par\smallskip
{\small Nota: \textit{gaps} médios com peso igual por instância. As duas falhas ILS permanecem na análise com a solução viável sem recursos. A coluna melhor média inclui empates; a última coluna conta somente vitórias exclusivas.\par}
{\fonte{elaboração própria a partir dos dados experimentais.}}
\end{singlespace}
\end{table}
